% Estado del proyecto LAAD de principio a fin: qué es, qué se hizo, cómo funciona, cómo se mide y qué falta.
% Compilar:  python scripts/cifras.py && python scripts/datos_presentacion.py && latexmk -pdf proyecto.tex
% Ninguna cifra se escribe a mano: las generales salen de ../articulo/cifras.tex (scripts/cifras.py) y las
% distribuciones y los resultados del harness de datos_proyecto.tex (scripts/datos_presentacion.py).
\documentclass[aspectratio=169,10pt]{beamer}
\usetheme{metropolis}
\metroset{block=fill,numbering=fraction,progressbar=frametitle,sectionpage=progressbar}
\usepackage[T1]{fontenc}
\usepackage[utf8]{inputenc}
\usepackage[spanish,es-noshorthands,es-tabla,es-nodecimaldot]{babel}
\usepackage{booktabs}
\usepackage{amssymb}
\usepackage{tikz}
\usetikzlibrary{positioning,arrows.meta,shapes.geometric,fit,calc,babel}
\usepackage{pgfplots}
\pgfplotsset{compat=1.18}
\usepackage{csquotes}
% GENERADO por scripts/cifras.py. No editar a mano: ejecutar el script.
% Lo cargan articulo/main.tex y presentacion/main.tex.
\newcommand{\cifraAfirmaciones}{504}
\newcommand{\cifraAlineaciones}{66}
\newcommand{\cifraAtributos}{30}
\newcommand{\cifraCandidatos}{19}
\newcommand{\cifraCatalogoRegistros}{66}
\newcommand{\cifraCitas}{1\,049}
\newcommand{\cifraClases}{41}
\newcommand{\cifraConceptos}{164}
\newcommand{\cifraDUnoEspurias}{285}
\newcommand{\cifraDSeisAfirmaciones}{504}
\newcommand{\cifraDSeisNegaciones}{20}
\newcommand{\cifraDefinicionesDeTrabajo}{22}
\newcommand{\cifraEntidadesNoPersona}{260}
\newcommand{\cifraEntidadesSinTipo}{135}
\newcommand{\cifraEsquemas}{12}
\newcommand{\cifraExpedientes}{1}
\newcommand{\cifraFormas}{50}
\newcommand{\cifraGrafos}{77}
\newcommand{\cifraHechosEnlazados}{13}
\newcommand{\cifraHechosNarrados}{61}
\newcommand{\cifraMenciones}{545}
\newcommand{\cifraMencionesEnlazadas}{123}
\newcommand{\cifraMutacionesDeclaradas}{36}
\newcommand{\cifraNegaciones}{20}
\newcommand{\cifraNegacionesConRespuesta}{4}
\newcommand{\cifraPiezas}{33}
\newcommand{\cifraPreguntas}{20}
\newcommand{\cifraPropiedades}{75}
\newcommand{\cifraRazonadorCierre}{189\,555}
\newcommand{\cifraRazonadorInconsistencias}{0}
\newcommand{\cifraRazonadorTripletas}{56\,351}
\newcommand{\cifraReconstrucciones}{35}
\newcommand{\cifraReconstruccionesConFechas}{36}
\newcommand{\cifraTerminosPropios}{0}
\newcommand{\cifraTripletas}{31\,037}
\newcommand{\cifraValoresLiterales}{211}

% Generado por scripts/datos_presentacion.py. No editar a mano.
\newcommand{\serieClases}{(0,53) (1,18) (2,16) (3,15) (4,13) (5,3)}
\newcommand{\rotulosClases}{{persona},{lugar},{cosa},{grupo},{hecho},{cargo}}
\newcommand{\serieCreencia}{(0,29) (1,70) (2,16) (3,2) (4,1)}
\newcommand{\rotulosCreencia}{{verdadero},{probable},{posible},{improbable},{indeterminado}}
\newcommand{\serieFuerza}{(0,249) (1,86) (2,53) (3,52) (4,43) (5,21)}
\newcommand{\rotulosFuerza}{{testimonial},{informativa},{resolutiva},{alegación de parte},{dictaminadora},{fedataria}}
\newcommand{\serieFuerzaX}{(249,0) (86,1) (53,2) (52,3) (43,4) (21,5)}
\newcommand{\datoAfirma}{484}
\newcommand{\datoNiega}{20}
\newcommand{\seriePiezas}{(0,36) (1,5) (2,30) (3,29) (4,19) (5,52) (6,26) (7,4) (8,16) (9,11) (10,5) (11,1) (12,7) (13,17) (14,7) (15,3) (16,22) (17,24) (18,19) (19,49) (20,6) (21,31) (22,6) (23,3) (24,42) (25,9) (26,17) (27,4) (28,4)}
\newcommand{\rotulosPiezas}{{P01},{P02},{P03},{P04},{P05},{P06},{P07},{P08},{P09},{P10},{P11},{P12},{P13},{P14},{P16},{P17},{P18},{P19},{P20},{P21},{P22},{P23},{P24},{P25},{P27},{P28},{P29},{P30},{P31}}
\newcommand{\nPiezasConAfirmaciones}{29}
\newcommand{\serieAgenteCompleta}{(0,3) (1,5) (2,5) (3,7) (4,8) (5,3) (6,0)}
\newcommand{\serieAgenteEsencial}{(0,3) (1,11) (2,12) (3,10) (4,12) (5,9) (6,0)}
\newcommand{\serieAgenteForma}{(0,14) (1,17) (2,19) (3,17) (4,18) (5,12) (6,0)}
\newcommand{\serieAgenteVerificada}{(0,19) (1,19) (2,18) (3,18) (4,19) (5,15) (6,0)}
\newcommand{\rotulosAgente}{{v1},{h1},{h2},{r1},{r2},{r3},{r4}}
\newcommand{\tablaAgente}{%
v1 & 20 & 3 & 3 & 14 & 19 & 25\,k \\
h1 & 20 & 5 & 11 & 17 & 19 & 36\,k \\
h2 & 20 & 5 & 12 & 19 & 18 & 31\,k \\
r1 & 20 & 7 & 10 & 17 & 18 & 31\,k \\
r2 & 20 & 8 & 12 & 18 & 19 & 33\,k \\
r3 & 15 & 3 & 9 & 12 & 15 & 27\,k \\
r4 & 0 & 0 & 0 & 0 & 0 & 0\,k \\
}
\newcommand{\mediaCompleta}{7,5}
\newcommand{\numCompleta}{7.50}
\newcommand{\desvCompleta}{0,5}
\newcommand{\mediaEsencial}{11,0}
\newcommand{\numEsencial}{11.00}
\newcommand{\desvEsencial}{1,0}
\newcommand{\mediaForma}{17,5}
\newcommand{\numForma}{17.50}
\newcommand{\desvForma}{0,5}
\newcommand{\mediaVerificada}{18,5}
\newcommand{\numVerificada}{18.50}
\newcommand{\desvVerificada}{0,5}
\newcommand{\nEjecucionesLlenas}{2}
\newcommand{\matrizEsencial}{(0,0) [1] (1,0) [0] (2,0) [0] (3,0) [0] (4,0) [0] (5,0) [0] (6,0) [0] (7,0) [0] (8,0) [0] (9,0) [0] (10,0) [0] (11,0) [0] (12,0) [0] (13,0) [0] (14,0) [1] (15,0) [0] (16,0) [0] (17,0) [0] (18,0) [0] (19,0) [1] (0,1) [1] (1,1) [0] (2,1) [1] (3,1) [1] (4,1) [1] (5,1) [0] (6,1) [0] (7,1) [1] (8,1) [1] (9,1) [0] (10,1) [0] (11,1) [1] (12,1) [0] (13,1) [1] (14,1) [1] (15,1) [0] (16,1) [1] (17,1) [0] (18,1) [1] (19,1) [0] (0,2) [1] (1,2) [0] (2,2) [1] (3,2) [1] (4,2) [1] (5,2) [0] (6,2) [0] (7,2) [0] (8,2) [1] (9,2) [1] (10,2) [1] (11,2) [0] (12,2) [0] (13,2) [1] (14,2) [1] (15,2) [0] (16,2) [1] (17,2) [0] (18,2) [1] (19,2) [1] (0,3) [1] (1,3) [0] (2,3) [1] (3,3) [1] (4,3) [1] (5,3) [0] (6,3) [0] (7,3) [0] (8,3) [1] (9,3) [0] (10,3) [0] (11,3) [0] (12,3) [0] (13,3) [1] (14,3) [1] (15,3) [0] (16,3) [0] (17,3) [1] (18,3) [1] (19,3) [1] (0,4) [1] (1,4) [0] (2,4) [1] (3,4) [1] (4,4) [1] (5,4) [0] (6,4) [0] (7,4) [0] (8,4) [1] (9,4) [0] (10,4) [1] (11,4) [1] (12,4) [0] (13,4) [1] (14,4) [1] (15,4) [0] (16,4) [0] (17,4) [1] (18,4) [1] (19,4) [1] (0,5) [1] (1,5) [0] (2,5) [1] (3,5) [1] (4,5) [1] (5,5) [0] (6,5) [0] (7,5) [0] (8,5) [1] (9,5) [0] (10,5) [1] (11,5) [1] (12,5) [0] (13,5) [1] (14,5) [1] (15,5) [-1] (16,5) [-1] (17,5) [-1] (18,5) [-1] (19,5) [-1] (0,6) [-1] (1,6) [-1] (2,6) [-1] (3,6) [-1] (4,6) [-1] (5,6) [-1] (6,6) [-1] (7,6) [-1] (8,6) [-1] (9,6) [-1] (10,6) [-1] (11,6) [-1] (12,6) [-1] (13,6) [-1] (14,6) [-1] (15,6) [-1] (16,6) [-1] (17,6) [-1] (18,6) [-1] (19,6) [-1]}
\newcommand{\maxEjecucion}{6}
\newcommand{\pcSiempreBien}{01, 03, 04, 05, 09, 14, 15, 19}
\newcommand{\pcSiempreMal}{02, 06, 07, 13, 16}
\newcommand{\pcVariables}{08, 10, 11, 12, 17, 18, 20}
\newcommand{\interCompleta}{10}
\newcommand{\interEsencial}{17}
\newcommand{\interForma}{20}
\newcommand{\interVerificada}{20}


\definecolor{tinta}{HTML}{23373B}
\definecolor{ocre}{HTML}{B5651D}
\definecolor{verde}{HTML}{3F7F4F}
\definecolor{rojo}{HTML}{B3261E}
\definecolor{ambar}{HTML}{C48A10}
\definecolor{gris}{HTML}{C9D1D2}
\setbeamercolor{alerted text}{fg=ocre}
\setbeamercolor{progress bar}{fg=ocre}
\newcommand{\cita}[1]{«\textit{#1}»}
\newcommand{\ttt}[1]{\texttt{\small #1}}
\newcommand{\cifra}[1]{\textbf{\alert{#1}}}
\newcommand{\si}{\textcolor{verde}{\checkmark}}
\newcommand{\no}{\textcolor{rojo}{$\times$}}
\newcommand{\aun}{\textcolor{ambar}{$\circ$}}
\tikzset{
  caja/.style={draw=tinta!60, fill=#1, rounded corners=3pt, align=center, font=\footnotesize, inner sep=4pt},
  caja/.default=white,
  flecha/.style={-{Stealth[length=2mm]}, thick, tinta!70},
  capa/.style={draw=tinta!60, fill=#1, rounded corners=3pt, minimum height=0.75cm, align=left, text width=12.8cm,
    font=\footnotesize},
}
\pgfplotsset{
  barras/.style={ybar, bar width=10pt, axis lines*=left, ymajorgrids, grid style={gris!60}, enlarge x limits=0.12,
    nodes near coords, every node near coord/.append style={font=\scriptsize}, tick label style={font=\scriptsize},
    ymin=0, xtick=data, height=4.6cm, width=7cm, axis line style={tinta!50}},
}

\title{LAAD: estado del proyecto}
\subtitle{Grafo de conocimiento de un expediente judicial colonial del ARC Cusco y su agente de consulta}
\author{Jhon Jesus Quispe Machaca \and Ian Logan Will Quispe Ventura}
\institute{Universidad Nacional de San Antonio Abad del Cusco}
\date{Repaso completo · \today}

\begin{document}
\maketitle

\begin{frame}{Qué contiene esta presentación}
\begin{columns}[T]
\column{0.48\textwidth}
\begin{enumerate}
\item El problema y el expediente
\item Materiales y restricciones
\item Cronología de lo hecho
\item El modelo: cuatro capas
\item Cómo se construye el grafo
\item El grafo en cifras
\end{enumerate}
\column{0.48\textwidth}
\begin{enumerate}\setcounter{enumi}{6}
\item Cómo se valida
\item Consulta y escala
\item El harness: agente de consulta
\item Resultados del harness
\item Problemas abiertos y propuestas
\item Pendientes, uso y glosario
\end{enumerate}
\end{columns}
\medskip
\footnotesize Toda cifra de estas diapositivas sale de los datos mediante un guion; ninguna se escribió a mano.
\end{frame}

% =====================================================================
\section{El problema y el expediente}

\begin{frame}{Qué es LAAD}
Un \alert{grafo de conocimiento} de expedientes judiciales coloniales del Fondo Intendencia (Archivo Regional del
Cusco), y un \alert{agente} que responde preguntas sobre él sin inventar.
\begin{itemize}
\item \textbf{Piloto:} Causa criminal, Leg.~14, Exp.~18 (Tinta, 1806). Manuel Huancachoque, indio, se queja del
      cacique recaudador Andrés Chuquitapa: lo acusa de haber sido coronel de Túpac Amaru y de abusos.
\item \textbf{Producto 1, el grafo:} RDF con estándares (RiC-O, CIDOC CRM, CRMinf, PiCo, ROAR, PROV-O, SKOS),
      perfil de aplicación \alert{OHAC}, sin términos propios.
\item \textbf{Producto 2, el harness:} herramientas con nombre sobre el grafo y un modelo de lenguaje que las usa;
      cada cita de la respuesta se comprueba contra lo que devolvieron.
\item \textbf{Producto 3, el artículo:} sobre la ontología (borrador en \ttt{articulo/}).
\end{itemize}
\end{frame}

\begin{frame}{Un expediente no registra hechos}
Registra lo que \alert{alguien afirmó}, ante quién, en qué actuación y con qué fuerza procesal.
\begin{columns}[T]
\column{0.48\textwidth}
\begin{block}{Agosto: certificación de Isidro Cano}
Chuquitapa fue \cita{complice y coronel} de Túpac Amaru.
\end{block}
\column{0.48\textwidth}
\begin{block}{Septiembre: declaración jurada del mismo Cano}
\cita{Comisionado por este y no coronel}. Otro testigo niega haber firmado; un tercero no sabe escribir.
\end{block}
\end{columns}
\medskip
El fiscal concluye que las quejas son \cita{menos conforme a la verdad}. ¿Quién tiene razón?
\alert{El grafo no lo decide}: guarda las voces y calcula dónde chocan.
\end{frame}

\begin{frame}{Por qué se rehízo: el prototipo anterior}
El primer intento (\ttt{../Grafo-Conocimiento-Cusco}, grafo de propiedades en Neo4j poblado por un modelo de lenguaje)
cometió tres errores de fondo:
\begin{enumerate}
\item Cada afirmación se volvió una tripleta \alert{sobre el mundo} (\enquote{Chuquitapa ES coronel}).
\item La resolución del tribunal se tomó como \alert{valor de verdad}.
\item Todas las grafías de un nombre se \alert{fusionaron} en una persona.
\end{enumerate}
Resultado: la frase del fiscal quedó atribuida a ocho declarantes, entre ellos el propio denunciante; la relación
\ttt{CONTRADICE} tenía 0 instancias.

\medskip
\textbf{Lección que guía todo LAAD:} el modelo de lenguaje propone; nunca decide identidades, verdades ni folios.
\end{frame}

% =====================================================================
\section{Materiales y restricciones}

\begin{frame}{Con qué se trabaja}
\begin{columns}[T]
\column{0.5\textwidth}
\textbf{Materiales}
\begin{itemize}\small
\item Transcripción del expediente (\ttt{corpus/}), identificada por su \ttt{sha256}.
\item Catálogo impreso del ARC (fichas e índices): capa C.
\item Especificaciones de los estándares copiadas en \ttt{vendor/} (49 archivos verificados).
\end{itemize}
\textbf{Hallazgos}
\begin{itemize}\small
\item La transcripción no marca recto ni vuelto; solo hay 4 marcas de folio (20, 28, 33, 34).
\item El \enquote{corpus modernizado} corrige en silencio (1809 → 1805, Mara → Mata): \alert{excluido}.
\item El final de UT20 es un documento republicano ajeno (P32).
\end{itemize}
\column{0.48\textwidth}
\textbf{Restricciones fijadas por el equipo}
\begin{itemize}\small
\item Sin imágenes de folios: el folio y el lado quedan \enquote{no consta}.
\item Un único validador humano (el autor); no hay acuerdo entre anotadores.
\item La unidad citable es el \alert{tramo} (UT00–UT20), no el folio.
\item Única normalización de citas: colapsar espacios. \enquote{Huancachoque} por \enquote{Guancachoque} se
      rechaza.
\item El portátil se calienta: nada de procesos en paralelo.
\end{itemize}
\end{columns}
\end{frame}

% =====================================================================
\section{Cronología de lo hecho}

\begin{frame}{Del 12 al 24 de noviembre de 2026}
\centering
\begin{tikzpicture}[x=1cm,y=1cm, every node/.style={font=\scriptsize}]
\draw[thick, tinta!60] (0,0) -- (8.3,0);
\draw[tinta!40, dashed, thick] (8.3,0) -- node[above, tinta!70] {pausa} (11.7,0);
\draw[thick, tinta!60] (11.7,0) -- (13.4,0);
\foreach \x/\d in {1.5/12 nov, 4.9/13 nov, 8.3/14 nov, 11.7/24 nov} {
  \fill[ocre] (\x,0) circle (3pt); \node[above=3pt] at (\x,0) {\textbf{\d}};
}
\foreach \x/\c/\t in {
  1.5/tinta!5/{Fuentes e índice; preguntas de competencia; diagnóstico del prototipo; modelo de 4 capas; revisiones 1 y 2; alineaciones; borrador del artículo},
  4.9/tinta!5/{Expediente completo anotado por agentes; tesauro 164/164 definido; capa 4 auditada; catálogo; 20/20 PC},
  8.3/tinta!5/{Varios expedientes; tesauro que crece; Oxigraph; cifras de una sola fuente; conectividad (35 $\to$ 118 reconstrucciones); presentación},
  11.7/ocre!12/{\textbf{Harness}: 14 herramientas, verificador, banco de 20 PC; cardinalidad; agente independiente del proveedor; Gemini; 6 bancos; decisiones del validador}} {
  \node[caja=\c, text width=2.9cm, minimum height=2.6cm, anchor=north] at (\x,-0.35) {\t};
}
\end{tikzpicture}

\medskip
\footnotesize Cada entregable vive en una rama git integrada en \ttt{main} con \ttt{merge --no-ff} (13 integraciones); el
harness está en \ttt{harness/consulta}, aún sin integrar. Detalle: \ttt{docs/bitacora.md} y \ttt{docs/decisiones/README.md}
(47 decisiones).
\end{frame}

% =====================================================================
\section{El modelo}

\begin{frame}{Cuatro capas, cada una en sus grafos nombrados}
\centering
\begin{tikzpicture}[node distance=1.4mm]
\node[capa=ocre!12] (c4) {\textbf{4 · Reconstrucciones} — identidades y hechos \emph{propuestos}, con autor, fecha, argumento y
  creencia; las hipótesis rivales coexisten (CRMinf, PiCo, ROAR)};
\node[capa=tinta!10, below=of c4] (c3) {\textbf{3 · Afirmaciones atribuidas} — quién afirma, sobre quién, qué atributo, qué valor,
  polaridad y fuerza (CIDOC CRM E13)};
\node[capa=tinta!6, below=of c3] (c2) {\textbf{2 · Menciones} — cada aparición escrita de un nombre, sin fusionar, con su cita literal
  (PiCo, ROAR, Web Annotation)};
\node[capa=tinta!3, below=of c2] (c1) {\textbf{1 · Documento} — expediente, piezas, actuaciones, tramos (RiC-O, CIDOC CRM)};
\node[capa=white, below=of c1] (p) {\textbf{P · Procedencia} — qué proceso, qué agente, qué versión de la transcripción (PROV-O);
  \textbf{C · Catálogo} — fichas del ARC como dato a contrastar};
\end{tikzpicture}

\smallskip
\footnotesize Sin valores de verdad en las capas 1--3. Solo la capa 4 tiene creencias. Las disputas no se guardan: se
\alert{calculan}.
\end{frame}

\begin{frame}{Cero términos propios: un perfil, no una ontología nueva}
OHAC declara \cifra{\cifraTerminosPropios} clases y propiedades. Se compone de:
\begin{itemize}
\item importaciones con versión fijada: RiC-O 1.1, CIDOC CRM 7.1.3, CRMinf 1.2.1, PiCo, ROAR, PROV-O, SKOS,
      Web Annotation;
\item un \alert{tesauro SKOS} de \cifra{\cifraEsquemas} esquemas y \cifra{\cifraConceptos} conceptos definidos
      (cargos, condiciones, tipos de actuación, roles, \cifra{\cifraAtributos} atributos, polaridad, fuerza…);
\item \cifra{\cifraFormas} \alert{formas SHACL} que hacen cumplir el perfil;
\item \cifra{\cifraAlineaciones} alineaciones externas, solo \ttt{closeMatch} o \ttt{relatedMatch}.
\end{itemize}
\medskip
Lo andino y lo procesal son \alert{datos}, no clases: el ayllu es un grupo, no un lugar; \cita{India Ordinaria} no se
iguala a \emph{originaria}. El tesauro crece con \alert{candidatos} (\cifra{\cifraCandidatos} hoy, con aviso SHACL) y
lleva versión semántica (0.2.0).
\end{frame}

\begin{frame}{Una afirmación en el grafo: la negación de Cano}
\small
\begin{tabular}{@{}ll@{}}
\toprule
Quién afirma (\ttt{P14}) & mención «Isidro Cano», declarante \\
En qué actuación (\ttt{P9i}) & declaración jurada de septiembre (pieza P18) \\
Sobre quién (\ttt{P140}) & mención «Andres Chuquitapa» (sujeto elíptico, marcado) \\
Qué atributo (\ttt{P177}) & \emph{ejerce-cargo} \\
Qué valor (\ttt{P141}) & \emph{coronel rebelde} \\
Polaridad y fuerza (\ttt{P2}) & \alert{niega} · testimonial \\
Responde a (\ttt{P17}) & la afirmación de agosto (A-UT01-001) \\
Cita literal & \cita{Comisionado por este y no coronel} (tramo UT11) \\
\bottomrule
\end{tabular}

\medskip
Ningún nodo dice si Chuquitapa fue coronel. La pregunta se responde como \alert{disputa}: quién lo afirma, quién lo
niega y qué autoridad se pronunció (ninguna).
\end{frame}

\begin{frame}{Identidades: hipótesis con autor y creencia}
\begin{columns}[T]
\column{0.55\textwidth}
\small
Las menciones no se fusionan. Una \alert{reconstrucción} de la capa 4 propone que varias menciones son la misma
persona, lugar o hecho, con:
\begin{itemize}
\item autor (hoy, un modelo de lenguaje) y fecha;
\item argumento escrito;
\item creencia: verdadero, probable, posible, improbable;
\item \alert{rivales}: lecturas incompatibles que comparten menciones.
\end{itemize}
Ejemplo: el \enquote{Mariano Casorla} que declara en P19, ¿es el granadero \enquote{Martin Cazorla}? Hay tres hipótesis y
el grafo guarda las tres.
\column{0.43\textwidth}
\begin{tikzpicture}[node distance=6mm and 4mm, every node/.style={font=\scriptsize}]
\node[caja=tinta!5] (m1) {«Martin Cazorla»\\P05};
\node[caja=tinta!5, right=of m1] (m2) {«Mariano Casorla»\\P19};
\node[caja=ocre!15, below=12mm of m1] (ra) {Hipótesis A\\misma persona\\\textit{probable}};
\node[caja=ocre!15, below=12mm of m2] (rb) {Hipótesis B\\otra persona\\\textit{improbable}};
\draw[flecha] (m1) -- (ra); \draw[flecha] (m2) -- (ra); \draw[flecha] (m2) -- (rb);
\draw[ocre, dashed, thick] (ra) -- node[below, font=\tiny] {rivales} (rb);
\end{tikzpicture}
\end{columns}
\end{frame}

\begin{frame}{Disputas calculadas y cardinalidad}
Una consulta compara las afirmaciones de un mismo sujeto y atributo:
\begin{itemize}
\item \textbf{se afirma y se niega} el mismo valor: disputa (coronel / no coronel);
\item \textbf{valores distintos}: disputa \emph{solo} si el atributo admite un valor (decisión 046, tabla aprobada por el
      validador): edad, fecha, duración, número de participantes, quién redactó… son \enquote{uno}; cargo, condición,
      parentesco… son \enquote{varios} (cacique y recaudador no se contradicen);
\item \textbf{identidad}: dos grafías comparan como iguales solo si pertenecen a una misma reconstrucción única; las
      menciones de reconstrucciones rivales nunca se unifican.
\end{itemize}
\medskip
La forma de cada respuesta también se calcula: \textbf{D} documental, \textbf{A} afirmaciones concordantes,
\textbf{DISPUTA}, \textbf{NC} no consta, \textbf{R} reconstrucción, \textbf{P} procedencia.
\end{frame}

% =====================================================================
\section{Cómo se construye el grafo}

\begin{frame}{De la transcripción al grafo}
\centering
\begin{tikzpicture}[node distance=5mm and 4mm, every node/.style={font=\scriptsize}]
\node[caja=tinta!5, text width=1.9cm] (t) {Transcripción\\y manifiesto\\del expediente};
\node[caja=ocre!12, text width=1.9cm, right=of t] (a) {Agentes LLM anotan\\\alert{registros planos}\\con cita obligatoria};
\node[caja=tinta!5, text width=1.9cm, right=of a] (v) {Validación del lote\\cita = subcadena\\ids, vocabulario};
\node[caja=tinta!5, text width=1.9cm, right=of v] (m) {Mapeo determinista\\(código, sin LLM)};
\node[caja=tinta!12, text width=1.9cm, right=of m] (g) {Grafo\\TriG + N-Quads};
\draw[flecha] (t) -- (a); \draw[flecha] (a) -- (v); \draw[flecha] (v) -- (m); \draw[flecha] (m) -- (g);
\node[caja=white, text width=2.2cm, below=16mm of g] (s) {SHACL\\(formas OHAC)};
\node[caja=white, text width=2.2cm, left=of s] (r) {Razonador\\OWL 2 RL};
\node[caja=white, text width=2.2cm, left=of r] (c) {\ttt{cifras.py}\\artículo, docs};
\node[caja=ocre!12, text width=2.4cm, left=of c] (p) {Propuestas de la\\capa 4 (agentes)\\+ auditoría};
\draw[flecha] (g) -- (s); \draw[flecha] (s) -- (r); \draw[flecha] (r) -- (c);
\draw[flecha] (p) |- (m);
\end{tikzpicture}

\medskip
\begin{itemize}\small
\item El modelo de lenguaje solo rellena un esquema plano y cerrado; \alert{el paso al grafo es código}.
\item Las piezas y su tipo de actuación los cura una persona (\ttt{datos/piezas.jsonl}).
\item Construir: \ttt{python scripts/construir.py}; un expediente nuevo: \ttt{scripts/nuevo\_expediente.py}.
\end{itemize}
\end{frame}

\begin{frame}{Pensado para un fondo, no para un expediente}
\begin{itemize}
\item Cada expediente se declara en un \alert{manifiesto}: signatura, piezas, lotes, formato de transcripción.
\item Las capas 1--3 viven en grafos del expediente (\ttt{g:<exp>/documento}, \ttt{g:<exp>/afirmaciones/…}).
\item La capa 4 es común: una reconstrucción puede unir menciones de \alert{expedientes distintos}.
\item Las consultas se restringen a un expediente o recorren el corpus.
\end{itemize}
\medskip
Hoy hay \cifra{\cifraExpedientes} expediente real; la generalización se probó con expedientes sintéticos (un pleito
civil no necesitó términos nuevos). El segundo expediente real (CO Cusco L70 E15, sorteado a ciegas) espera su
transcripción.
\end{frame}

% =====================================================================
\section{El grafo en cifras}

\begin{frame}{Tamaño del piloto}
\begin{columns}[T]
\column{0.5\textwidth}
\begin{itemize}
\item \cifra{\cifraPiezas} piezas del expediente
\item \cifra{\cifraMenciones} menciones (\cifra{\cifraPersonas} de persona)
\item \cifra{\cifraAfirmaciones} afirmaciones (\cifra{\cifraNegaciones} negaciones)
\item \cifra{\cifraCitas} citas literales, todas subcadena exacta
\item \cifra{\cifraReconstrucciones} reconstrucciones
\item \cifra{\cifraCatalogoRegistros} fichas e índices del catálogo
\end{itemize}
\column{0.5\textwidth}
\begin{itemize}
\item \cifra{\cifraTripletas} tripletas en \cifra{\cifraGrafos} grafos nombrados
\item \cifra{\cifraClases} clases y \cifra{\cifraPropiedades} propiedades, todas de estándares
\item \cifra{\cifraConceptos} conceptos del tesauro (\cifra{\cifraDefinicionesDeTrabajo} con definición de trabajo)
\item \cifra{\cifraPreguntas} preguntas de competencia
\end{itemize}
\end{columns}
\end{frame}

\begin{frame}{La capa 3: quién habla y con qué fuerza}
\begin{columns}[T]
\column{0.5\textwidth}
\begin{tikzpicture}
\begin{axis}[xbar, bar width=8pt, xmin=0, ytick={0,...,5}, yticklabels/.expanded={\rotulosFuerza},
  height=5.2cm, width=6.4cm, enlarge y limits=0.12, enlarge x limits={upper,value=0.18}, xmajorgrids,
  grid style={gris!60}, y dir=reverse, title={\footnotesize Afirmaciones por fuerza institucional}, axis lines*=left,
  tick label style={font=\scriptsize}, nodes near coords, every node near coord/.append style={font=\scriptsize},
  axis line style={tinta!50}]
\addplot[fill=tinta!70, draw=none] coordinates {\serieFuerzaX};
\end{axis}
\end{tikzpicture}
\column{0.5\textwidth}
\small
\begin{itemize}
\item La mayoría son \alert{testimonios}: lo que declaran testigos y certificantes.
\item Las \alert{resolutivas} (autos, decretos) también son afirmaciones, no hechos: una orden no implica su
      cumplimiento.
\item \cifra{\datoNiega} de las \cifra{\cifraAfirmaciones} afirmaciones son negaciones, y
      \cifra{\cifraNegacionesConRespuesta} de ellas apuntan explícitamente a lo que niegan.
\end{itemize}
\end{columns}
\end{frame}

\begin{frame}{Afirmaciones a lo largo del expediente}
\centering
\begin{tikzpicture}
\begin{axis}[barras, bar width=7pt, width=14cm, height=5cm, xticklabels/.expanded={\rotulosPiezas},
  x tick label style={rotate=60, font=\tiny}, enlarge x limits=0.03, nodes near coords={},
  title={\footnotesize Afirmaciones por pieza (\nPiezasConAfirmaciones{} piezas con alguna)}]
\addplot[fill=ocre!80, draw=none] coordinates {\seriePiezas};
\end{axis}
\end{tikzpicture}

\footnotesize Concentran más voces P06 (escrito de Huancachoque), P21 (declaración de Orue) y P27 (informe del subdelegado Mata);
las piezas que faltan no contienen afirmaciones (trámite o documento ajeno).
\end{frame}

\begin{frame}{La capa 4: qué se ha reconstruido y con qué creencia}
\begin{columns}[T]
\column{0.5\textwidth}
\begin{tikzpicture}
\begin{axis}[barras, xticklabels/.expanded={\rotulosClases}, title={\footnotesize Reconstrucciones por clase}, width=6.6cm]
\addplot[fill=tinta!70, draw=none] coordinates {\serieClases};
\end{axis}
\end{tikzpicture}
\column{0.5\textwidth}
\begin{tikzpicture}
\begin{axis}[barras, xticklabels/.expanded={\rotulosCreencia}, x tick label style={rotate=25, font=\tiny},
  title={\footnotesize Creencia declarada}, width=6.6cm]
\addplot[fill=ocre!80, draw=none] coordinates {\serieCreencia};
\end{axis}
\end{tikzpicture}
\end{columns}
\footnotesize La creencia es la de quien propuso la reconstrucción (un modelo de lenguaje), no un hecho.
\cifra{\cifraReconstruccionesSinAuditar} de \cifra{\cifraReconstrucciones} no han pasado la auditoría del segundo modelo y
\alert{ninguna tiene validación humana}.
\end{frame}

\begin{frame}{Conectividad: qué parte del texto está enlazada}
\centering
\begin{tikzpicture}
\begin{axis}[xbar, bar width=9pt, xmin=0, xmax=100, width=11cm, height=4.4cm, axis lines*=left, xmajorgrids,
  grid style={gris!60}, ytick={0,1,2,3}, yticklabels={menciones, menciones de persona, hechos narrados, menciones no persona con tipo},
  y dir=reverse, enlarge y limits=0.2, tick label style={font=\scriptsize}, xlabel={\scriptsize \% enlazado},
  nodes near coords={\pgfmathprintnumber[fixed,precision=0]{\pgfplotspointmeta}\,\%}, every node near coord/.append style={font=\scriptsize}]
\addplot[fill=tinta!70, draw=none] coordinates {
  ({100*\cifraMencionesEnlazadas/\cifraMenciones},0) ({100*\cifraPersonasEnlazadas/\cifraPersonas},1)
  ({100*\cifraHechosEnlazados/\cifraHechosNarrados},2)
  ({100*(\cifraEntidadesNoPersona-\cifraEntidadesSinTipo)/\cifraEntidadesNoPersona},3)};
\end{axis}
\end{tikzpicture}

\footnotesize \cifraMencionesEnlazadas{} de \cifraMenciones{} menciones pertenecen a alguna reconstrucción;
\cifraHechosEnlazados{} de \cifraHechosNarrados{} hechos narrados están reconstruidos; \cifraEntidadesSinTipo{} menciones
que no son de persona carecen de tipo de entidad.
\end{frame}

\begin{frame}{Ver el grafo: el visor}
\begin{columns}[T]
\column{0.55\textwidth}
\small
\ttt{python scripts/visor.py} genera \ttt{visor/datos.js}; se abre \ttt{visor/index.html} en el navegador (necesita
conexión para la biblioteca de la red).
\begin{itemize}
\item \textbf{Nodos}: reconstrucciones (color por clase, borde por creencia), piezas, menciones sin identidad única.
\item \textbf{Aristas}: \enquote{afirma sobre} (rojo si niega, ámbar si está en disputa), hipótesis rivales
      (discontinua), valor.
\item \textbf{Ficha}: al pulsar un nodo, sus menciones y todas sus afirmaciones con cita, pieza, tramo y fuerza.
\item Búsqueda por nombre, grafía o id; filtro por pieza; \enquote{solo disputas}.
\end{itemize}
\column{0.42\textwidth}
\begin{block}{La misma regla que el harness}
Una mención se dibuja como su reconstrucción solo si pertenece a \alert{una}. Si pertenece a varias, queda como nodo
propio unido a cada hipótesis: el visor no elige.
\end{block}
\end{columns}
\end{frame}

% =====================================================================
\section{Cómo se valida}

\begin{frame}{Evaluación sin anotadores externos}
\small
\begin{tabular}{@{}lr@{}}
\toprule
Criterio & Resultado \\
\midrule
Citas que son subcadena exacta del tramo & \cifraCitas/\cifraCitas \\
Violaciones SHACL del grafo & 0 \\
Mutaciones del grafo que SHACL rechaza (pruebas negativas) & \cifraMutacionesDeclaradas/\cifraMutacionesDeclaradas \\
Inconsistencias bajo OWL 2 RL (\cifraRazonadorCierre{} tripletas tras el cierre) & \cifraRazonadorInconsistencias \\
Términos nuevos para un pleito civil sintético & 0 \\
Preguntas de competencia respondidas por consultas con nombre & \cifraPreguntas/\cifraPreguntas \\
\bottomrule
\end{tabular}

\medskip
Las 20 preguntas se escribieron \alert{antes} de modelar (\ttt{docs/B\_preguntas\_competencia.md}); cada una es una
prueba automática. Hay 21 archivos de pruebas en \ttt{tests/} (la suite completa tarda de 7 a 28 minutos).
\end{frame}

\begin{frame}{Lo que los estándares no esperaban}
Siete defectos hallados al validar y razonar con datos reales; no se corrigen en origen, se evitan o se miden:
\begin{itemize}\small
\item \textbf{D1} ROAR: un dominio global vuelve inconsistente el patrón de PiCo (\cifraDUnoEspurias{} tipificaciones
      espurias).
\item \textbf{D6} CRMinf 1.2.1: toda E13 es una argumentación \enquote{honesta}; cada testimonio hereda una
      \alert{presunción de sinceridad} (\cifraDSeisAfirmaciones{} afirmaciones, \cifraDSeisNegaciones{} negaciones).
\item \textbf{D7} CRM + PiCo: toda mención que afirma se infiere actor.
\item D2--D5: prefijos de PiCo, relaciones temporales heredadas, creencias literales, folios en RiC-O.
\end{itemize}
\medskip
\textbf{Revisiones}: dos revisiones \enquote{independientes} del modelo y del artículo, hechas por otras instancias del
modelo de lenguaje sin acceso a la justificación del autor; 3 hallazgos críticos y 8 importantes, atendidos.
\end{frame}

% =====================================================================
\section{Consulta y escala}

\begin{frame}{Consulta y escala}
\begin{itemize}
\item \cifra{16} consultas SPARQL con nombre en \ttt{consultas/}; el agente nunca escribe SPARQL.
\item Motor: \alert{Oxigraph} (rdflib queda como alternativa y una prueba exige las mismas filas en ambos).
      Una consulta de afirmaciones pasó de unos 15\,s a centésimas.
\item Escala, con el piloto replicado hasta 1\,000 veces: carga y espacio lineales; las consultas restringidas a un
      expediente tardan lo mismo con 1 que con 100; las que recorren el corpus (disputas) crecen con él.
\end{itemize}
\medskip
\footnotesize Mediciones con réplicas del piloto, no con el archivo real: indican cómo crece el coste
(\ttt{datos/grafo/escala.md}).
\end{frame}

% =====================================================================
\section{El harness: agente de consulta}

\begin{frame}{Arquitectura del harness}
\centering
\begin{tikzpicture}[node distance=5mm and 7mm, every node/.style={font=\scriptsize}]
\node[caja=tinta!5, text width=2.1cm] (q) {Pregunta en\\lenguaje natural};
\node[caja=ocre!15, text width=2.5cm, right=of q] (llm) {\textbf{Modelo de lenguaje}\\(proveedor intercambiable:\\hoy Gemini)};
\node[caja=tinta!10, text width=2.6cm, right=of llm] (h) {\textbf{14 herramientas}\\con nombre, JSON Schema\\e ids validados};
\node[caja=tinta!5, text width=2.2cm, right=of h] (k) {Consultas con nombre\\sobre Oxigraph};
\draw[flecha] (q) -- (llm); \draw[flecha, transform canvas={yshift=2mm}] (llm) -- node[above, font=\tiny] {llamada} (h);
\draw[flecha, transform canvas={yshift=-2mm}] (h) -- node[below, font=\tiny] {filas + texto} (llm);
\draw[flecha] (h) -- (k);
\node[caja=verde!12, text width=3.4cm, below=10mm of llm] (v) {\textbf{Verificador}: cada «cita» debe ser subcadena de
  una fila devuelta; cada id de afirmación, estar en ellas};
\node[caja=white, text width=2.6cm, right=of v] (r) {Respuesta + \ttt{FORMA:}\\verificada o\\NO VERIFICADA};
\draw[flecha] (llm) -- node[left, font=\tiny] {respuesta} (v); \draw[flecha] (v) -- (r);
\draw[flecha, dashed] (v.west) -| node[pos=0.25, below, font=\tiny] {1 corrección} ($(q.south)+(0.4,0)$) -- ($(llm.south west)+(0.2,0)$);
\end{tikzpicture}

\medskip
\begin{itemize}\footnotesize
\item El bucle, las instrucciones (\ttt{prompts/agente\_consulta\_v2.md}) y la verificación viven en LAAD; cada proveedor
      es un adaptador pequeño (\ttt{src/laad/proveedores/}). Cambiar de modelo es escribir un adaptador.
\item Los nombres escritos se resuelven con \ttt{buscar} (variantes gráficas: qu/c, hu/gu, b/v, y/i…), que devuelve
      \alert{todos} los candidatos con su creencia y sus rivales.
\end{itemize}
\end{frame}

\begin{frame}{Las 14 herramientas}
\scriptsize
\begin{tabular}{@{}lp{10.2cm}@{}}
\toprule
\ttt{buscar} & nombre escrito con cualquier grafía → ids de reconstrucciones y menciones \\
\ttt{afirmaciones} & quién dijo qué, con cita, tramo, rol y fuerza; calcula disputas y comprueba el dominio \\
\ttt{disputas} & disputas calculadas en el corpus para una entidad o un atributo \\
\ttt{disputas\_firma} & firmas que constan en el texto y que alguien niega haber puesto \\
\ttt{menciones\_de} & todas las grafías que agrupa una reconstrucción \\
\ttt{reconstrucciones} & hipótesis de la capa 4 con autor, creencia, argumento y rivales \\
\ttt{cargos} & ejercicios de cargo con vigencia \\
\ttt{piezas} · \ttt{control\_secuencia} & piezas en orden físico; piezas cuyo orden contradice sus fechas \\
\ttt{participantes} & quién actúa en cada pieza y con qué papel procesal \\
\ttt{procedencia} & transcripción (sha256), actividad, agente y validación de un dato \\
\ttt{catalogo} & fichas del catálogo del ARC pertinentes (todo es \enquote{según el catálogo}) \\
\ttt{identificadores\_externos} & lugares y grupos con identificador HGIS o WHG \\
\ttt{vocabulario} & lista cerrada de conceptos de un esquema del tesauro \\
\bottomrule
\end{tabular}

\medskip
Uso a mano, sin modelo: \ttt{python scripts/preguntar.py afirmaciones --entidad R-chuquitapa}.
\end{frame}

\begin{frame}{Cómo se mide: el banco de 20 preguntas}
\begin{itemize}\small
\item Cada PC tiene una \alert{referencia} calculada (\ttt{laad.preguntas}) y una forma esperada (D, A, DISPUTA, NC, R, P).
\item \textbf{Completa}: la respuesta trae todas las afirmaciones, afirmantes, reconstrucciones, citas e ids de la
      referencia, con la forma correcta y verificada. Es exigente: todo o nada.
\item \textbf{Esencial}: como completa, pero perdona citas e ids \alert{periféricos} (frases secundarias de la referencia).
\item \textbf{Forma} correcta y respuesta \textbf{verificada} (sin citas inventadas).
\item Los errores de herramienta que el agente corrige solo son \enquote{avisos}, no fallos.
\end{itemize}
\medskip
\small
Tres niveles medidos con el mismo comprobador:
\begin{enumerate}\small
\item planes escritos a mano (techo de las herramientas): \cifra{20/20};
\item un intermediario Claude que solo usa \ttt{preguntar.py} (techo práctico con un modelo fuerte);
\item el agente real con Gemini 3.5 Flash-Lite (nivel gratuito).
\end{enumerate}
\end{frame}

\begin{frame}{Decisiones del validador sobre las referencias (24-nov)}
\small
\begin{tabular}{@{}llp{8.2cm}@{}}
\toprule
PC & Decisión & Efecto \\
\midrule
PC02 & aprobada & la remisión \enquote{De folio 4} se responde con la afirmación de folio y los firmantes de la pieza \\
PC08 & forma \textbf{DISPUTA} & Cano se contradice (coronel; fecha del hecho); antes era A \\
PC13 & rechazada & se mantiene la referencia amplia: el fallo del agente es esperado \\
PC16 & P28 + P29 & la resolución son las 26 afirmaciones de esas dos piezas \\
\bottomrule
\end{tabular}

\medskip
Con ello, en la evaluación del artículo el criterio \enquote{disputa o no consta} pasa de 6/6 a \alert{7/7} (PC08 se suma).
\end{frame}

% =====================================================================
\section{Resultados del harness}

\begin{frame}{Evolución del agente (Gemini 3.5 Flash-Lite)}
\begin{columns}[T]
\column{0.56\textwidth}
\begin{tikzpicture}
\begin{axis}[width=8.4cm, height=5.6cm, xtick=data, xticklabels/.expanded={\rotulosAgente}, ymin=0, ymax=20.5,
  ytick={0,5,10,15,20}, ymajorgrids, grid style={gris!60}, axis lines*=left, tick label style={font=\scriptsize},
  legend style={font=\tiny, at={(0.5,-0.18)}, anchor=north, legend columns=4, draw=none}, ylabel={\scriptsize PC de 20},
  mark size=1.8pt]
\addplot[thick, verde, mark=*] coordinates {\serieAgenteVerificada};
\addplot[thick, tinta, mark=square*] coordinates {\serieAgenteForma};
\addplot[thick, ocre, mark=triangle*] coordinates {\serieAgenteEsencial};
\addplot[thick, rojo, mark=diamond*] coordinates {\serieAgenteCompleta};
\legend{verificadas, forma correcta, esencial, completas}
\end{axis}
\end{tikzpicture}
\column{0.42\textwidth}
\scriptsize
\begin{tabular}{@{}lrrrrrr@{}}
\toprule
& PC & compl. & esenc. & forma & verif. & tokens \\
\midrule
\tablaAgente
\bottomrule
\end{tabular}

\medskip
\footnotesize v1: instrucciones v1. h1, h2: instrucciones v2 y herramientas mejoradas. r1--r3: tres repeticiones de la
misma configuración (r3 se cortó en PC15 por la cuota diaria).
\end{columns}
\end{frame}

\begin{frame}{Media de las repeticiones y techo}
\begin{columns}[T]
\column{0.55\textwidth}
\begin{tikzpicture}
\begin{axis}[ybar=1pt, bar width=8pt, width=8cm, height=5.4cm, ymin=0, ymax=21, axis lines*=left, ymajorgrids,
  grid style={gris!60}, xtick={0,1,2,3}, xticklabels={completa, esencial, forma, verificada}, enlarge x limits=0.18,
  tick label style={font=\scriptsize}, legend style={font=\tiny, at={(0.5,-0.15)}, anchor=north, legend columns=3,
  draw=none}, area legend, nodes near coords, every node near coord/.append style={font=\tiny}]
\addplot[fill=tinta!25, draw=none] coordinates {(0,20) (1,20) (2,20) (3,20)};
\addplot[fill=ocre!70, draw=none] coordinates {(0,\interCompleta) (1,\interEsencial) (2,\interForma) (3,\interVerificada)};
\addplot[fill=tinta!80, draw=none] coordinates {(0,\numCompleta) (1,\numEsencial) (2,\numForma) (3,\numVerificada)};
\legend{planes a mano, intermediario Claude, Gemini lite (media)}
\end{axis}
\end{tikzpicture}
\column{0.43\textwidth}
\small
Gemini 3.5 Flash-Lite, media de \nEjecucionesLlenas{} repeticiones completas:
\begin{itemize}\footnotesize
\item completas \cifra{\mediaCompleta} ± \desvCompleta
\item esencial \cifra{\mediaEsencial} ± \desvEsencial
\item forma \cifra{\mediaForma} ± \desvForma
\item verificadas \cifra{\mediaVerificada} ± \desvVerificada
\end{itemize}
\footnotesize La dispersión es pequeña: dos o tres repeticiones bastan. El intermediario Claude es una cota superior
(medido hoy con las referencias vigentes); vio los diagnósticos, así que está algo \alert{contaminado}.
\end{columns}
\end{frame}

\begin{frame}{Qué preguntas resuelve, ejecución a ejecución}
\centering
\begin{tikzpicture}
\begin{axis}[width=14.5cm, height=5.2cm, xmin=-0.5, xmax=19.5, ymin=-0.5, ymax=5.5, y dir=reverse,
  xtick={0,...,19}, xticklabels={01,02,03,04,05,06,07,08,09,10,11,12,13,14,15,16,17,18,19,20},
  ytick={0,...,5}, yticklabels/.expanded={\rotulosAgente}, tick label style={font=\scriptsize}, xlabel={\scriptsize PC},
  axis line style={draw=none}, tick style={draw=none}, colormap={estado}{color(0)=(gris) color(1)=(rojo!55) color(2)=(verde!75)},
  point meta min=-1, point meta max=1]
\addplot[scatter, only marks, mark=square*, mark size=9.5pt, point meta=explicit, scatter/use mapped color={draw=white, fill=mapped color}]
  coordinates {\matrizEsencial};
\end{axis}
\end{tikzpicture}

\footnotesize Verde: lo esencial correcto. Rojo: no. Gris: sin medir. Siempre bien desde h1: PC01, 03, 04, 05, 09, 14, 15, 19.
Siempre mal: PC02, 06, 07, 13, 16. Variables: PC08, 10, 11, 12, 17, 18, 20.
\end{frame}

\begin{frame}{Cuánto cuesta}
\begin{itemize}
\item Unos \alert{30\,000 tokens por pregunta}: cada turno reenvía todo el historial (instrucciones, herramientas y filas).
\item Nivel gratuito de Google AI Studio, por modelo y día (se renueva a medianoche del Pacífico):
  \begin{itemize}\small
  \item \textbf{gemini-3.6-flash}: 20 peticiones/día, menos de un banco; inservible para medir;
  \item \textbf{gemini-3.5-flash-lite}: unas 500 peticiones/día, de 4 a 5 bancos.
  \end{itemize}
\item Un banco tarda unos 15 minutos (red, no CPU: el portátil se mantiene a 43--55\,°C).
\item \alert{Privacidad}: en el nivel gratuito Google puede usar lo enviado (preguntas y texto del expediente, que es
      público) para mejorar sus productos. Nunca se envían datos personales; la clave vive en un \ttt{.env} ignorado
      por git.
\end{itemize}
\end{frame}

% =====================================================================
\section{Problemas abiertos y propuestas}

\begin{frame}{Por qué fallan las que fallan}
\scriptsize
\begin{tabular}{@{}lp{3.9cm}p{4.2cm}l@{}}
\toprule
PC & Qué pasa & Causa & Estado \\
\midrule
PC02 & 0\,\% de las afirmaciones & lee \enquote{folio 4} como la pieza P04 & \aun{} corregido, sin medir \\
PC10 & forma D en vez de DISPUTA; faltan 15 afirmaciones & usa \ttt{participantes} en vez de \ttt{afirmaciones(redacto)} & \aun{} corregido, sin medir \\
PC08 & todo encontrado, respuesta \alert{no verificada} & el modelo lite retoca o une citas; el verificador lo rechaza & \no \\
PC16 & 15\,\% de las afirmaciones & no sabe que la resolución está en P28/P29: busca por persona & \no \\
PC07 & 18\,\% de los lugares & lista larga: demasiadas llamadas & \no \\
PC06 & faltan 7 de 19 reconstrucciones & lista larga; se le escapan las rivales & \no \\
PC12, PC17 & a veces, forma equivocada & confunde \enquote{un punto disputado} con \enquote{respuesta en disputa} & \aun{} parcial \\
PC13 & NC sobre una referencia amplia & decisión del validador & esperado \\
\bottomrule
\end{tabular}

\medskip
\footnotesize PC03, 04, 09, 11 y 14 están bien en lo esencial: solo les faltan citas secundarias de la referencia. Por eso
\enquote{esencial} describe mejor el agente que \enquote{completa}.
\end{frame}

\begin{frame}{Qué se propone, por orden de rendimiento esperado}
\begin{enumerate}\small
\item \textbf{Medir PC02 y PC10} con la cuota renovada (dos repeticiones nuevas). Si salen, \enquote{esencial} sube a unas 13.
\item \textbf{Citar por id}: el modelo escribe \ttt{[A-UT11-004]} y el sistema inserta la cita literal del grafo. Las
      citas inventadas dejan de ser posibles \alert{por construcción}. Resuelve PC08.
\item \textbf{Listas completas en una llamada}: todas las reconstrucciones de una clase, compactas, con identificador HGIS
      y rivales; y piezas con su tipo de acto para ir a la resolución. Para PC06, PC07, PC16.
\item \textbf{Avisar del desacuerdo de forma} en la vuelta de corrección, sin imponerla. Para PC12, PC17.
\item \textbf{Informar \enquote{esencial} como cifra principal}, con PC13 aparte.
\item Menor: \ttt{buscar} mezcla clases, \ttt{catalogo} devuelve demasiado, \ttt{valor} busca por subcadena.
\end{enumerate}
\medskip
\small
\alert{El límite está en las herramientas más que en el modelo}: Claude, con las mismas herramientas, se queda en
\interEsencial{} de 20 en lo esencial. Por eso 2 y 3 van antes que cambiar de modelo.
\end{frame}

% =====================================================================
\section{Pendientes, uso y glosario}

\begin{frame}{Qué falta}
\begin{columns}[T]
\column{0.5\textwidth}
\textbf{Producto (prioridad)}
\begin{itemize}\footnotesize
\item Harness: puntos 1--5 de la diapositiva anterior.
\item Integrar \ttt{harness/consulta} en \ttt{main} cuando el resultado sea bueno.
\item Documentar el harness de principio a fin (bitácora, D §4.2, decisiones, artículo, preguntas de competencia).
\end{itemize}
\textbf{Validación humana (solo el autor)}
\begin{itemize}\footnotesize
\item Muestra de unas 15 afirmaciones por registro.
\item \cifraCandidatos{} conceptos candidatos por aprobar o rechazar.
\item \cifraDefinicionesDeTrabajo{} definiciones de trabajo sin fuente.
\item \cifraReconstruccionesSinAuditar{} reconstrucciones sin auditar.
\end{itemize}
\column{0.48\textwidth}
\textbf{Datos}
\begin{itemize}\footnotesize
\item Segundo expediente real: falta su transcripción.
\item Renombrar el lote \ttt{anotacion-manual-expediente-v1} (lo anotaron agentes, no a mano).
\item Problemas de anotación de la capa 3 ya listados (\ttt{hechos\_observaciones.md}).
\end{itemize}
\textbf{Aparcado por decisión del equipo}
\begin{itemize}\footnotesize
\item Licencia, DOI (Zenodo), w3id, declaración de uso de IA.
\item Sede del artículo (ESWC 2027, Resources Track) y recorte a su extensión.
\end{itemize}
\end{columns}
\end{frame}

\begin{frame}{Límites que hay que decir siempre}
\begin{itemize}
\item Un solo expediente real ($n=1$).
\item Anotaciones y reconstrucciones propuestas por modelos de lenguaje; un único validador humano, sin acuerdo
      entre anotadores; las \enquote{revisiones independientes} fueron otras instancias del modelo.
\item Sin imágenes: el folio y el lado de cada tramo no constan, salvo cuatro marcas del texto.
\item OWL 2 RL es incompleto: la ausencia de inconsistencias es necesaria, no suficiente.
\item El harness se mide con un banco de 20 preguntas escritas por los autores: sirve para comparar versiones, no
      para afirmar qué tan bien responde a cualquier usuario.
\end{itemize}
\end{frame}

\begin{frame}{Cómo se usa}
\scriptsize
\begin{tabular}{@{}lp{10.4cm}@{}}
\toprule
Construir el grafo & \ttt{python scripts/construir.py} \\
Cifras & \ttt{scripts/cifras.py}, \ttt{scripts/datos\_presentacion.py} \\
Ver el grafo & \ttt{python scripts/visor.py}, luego abrir \ttt{visor/index.html} \\
Consultar a mano & \ttt{python scripts/preguntar.py buscar --nombre Chuquitapa} \\
Preguntar al agente & \ttt{python scripts/agente.py --env <archivo .env> "pregunta"} \\
Banco del agente & \ttt{python scripts/agente.py --banco --sufijo r4} \\
Comparar bancos & \ttt{python scripts/comparar\_bancos.py datos/evaluacion/agente\_*.jsonl} \\
Planes a mano & \ttt{python scripts/evaluar\_harness.py} \\
Pruebas & \ttt{pytest tests/test\_herramientas.py} (la suite completa, solo al final) \\
\bottomrule
\end{tabular}

\medskip
\textbf{Mapa}: \ttt{ontologia/} perfil y tesauro · \ttt{datos/} registros, grafo, evaluación · \ttt{consultas/} SPARQL ·
\ttt{src/laad/} código · \ttt{prompts/} instrucciones del agente · \ttt{docs/} documentación (índice:
\ttt{docs/00\_indice.md}) · \ttt{articulo/} · \ttt{presentacion/} · \ttt{visor/}.
\end{frame}

\begin{frame}{Glosario mínimo}
\footnotesize
\begin{tabular}{@{}lp{10.6cm}@{}}
\toprule
RDF, TriG, N-Quads & formato de grafo en tripletas; TriG y N-Quads admiten \alert{grafos nombrados} (uno por capa y expediente) \\
E13 (CIDOC CRM) & \enquote{asignación de atributo}: alguien atribuye a algo un valor; base de la capa 3 \\
PiCo, ROAR & modelos de personas \emph{tal como constan} en las fuentes, y su reconstrucción \\
CRMinf & creencias e inferencias: la capa 4 dice quién cree qué y con qué grado \\
SKOS, tesauro & vocabulario controlado de conceptos (cargos, roles, atributos…) \\
SHACL, forma & reglas que el grafo debe cumplir; una violación es un error, un aviso no \\
OWL 2 RL & razonador que detecta contradicciones lógicas entre los estándares y los datos \\
SPARQL, Oxigraph & lenguaje de consulta de grafos RDF y el motor que lo ejecuta \\
PC & pregunta de competencia: lo que el grafo debe poder responder \\
Harness & herramientas + modelo de lenguaje + verificador + banco de medición \\
Tramo (UT) & unidad de la transcripción que se cita; reemplaza al folio \\
\bottomrule
\end{tabular}
\end{frame}

\begin{frame}[standout]
Estado: el grafo está completo y validado para un expediente.\\[4pt]
El agente acierta lo esencial en 11 de 20 preguntas;\\ las herramientas permiten 20 de 20.\\[4pt]
\alert{Siguiente: citar por id y listas completas.}
\end{frame}

\end{document}
